\setupcolors[state=start]
\setupcolor[hex]
\definecolor[darkblue][h=003366]
\definecolor[darkgray][h=e6e6e6]
\definecolor[redline][h=d9534f]
\definecolor[blueline][h=5bc0de]
\definecolor[greenline][h=5cb85c]

\setuphead[section][style=\bfa, color=darkblue, before=\blank, after=\blank]
\setuphead[subsection][style=\bf, color=black, before=\blank, after=\blank]
\setuppapersize[A4][A4]
\setupbodyfont[11pt]
\setuplayout[width=middle, height=middle, topspace=2cm, bottomspace=2cm]

\starttext

\title{\bf Análisis Comparativo Avanzado y Escalabilidad Paramétrica en Modelos de Linear Picking Routing}
\blank[big]

Este documento presenta un análisis empírico y teórico exhaustivo sobre el desempeño operativo, la eficiencia computacional y las fronteras de escalabilidad de las tres arquitecturas de optimización implementadas para resolver el problema integrado de preparación de pedidos y ruteo en almacenes funcionales ({\em Linear Picking Routing Problem}, LPRP): el modelo exacto de {\bf Programación Lineal Entera Mixta (MILP - OPP\_v3)} resuelto mediante SCIP, la heurística de búsqueda local {\bf Hill Climbing (HC)} y la metaheurística bio-inspirada {\bf Ant Colony Optimization (ACO)}. 

Los datos analizados provienen de una muestra de más de 200 experimentos controlados extraídos del consolidado final de pruebas de estrés masivas (\type{complete_output_v3_stress_tests.csv}) y pruebas normales, evaluando la escalabilidad paramétrica frente a variables críticas como el volumen de órdenes, la dimensión topológica del layout y restricciones de capacidad vehicular.

\startsection[title={1. Estadísticas Descriptivas Generales y Metodología de Agregación}]

El universo de datos crudos fue procesado mediante herramientas de manipulación de matrices lógicas bajo un estricto protocolo de depuración estadística. La evaluación de las métricas temporales (Media, Desvío Estándar y Máximo) se restringió exclusivamente a aquellas instancias experimentales que alcanzaron una convergencia exitosa (estado {\em ok}). Esto previene el sesgo matemático que introduciría la inclusión de escenarios truncados por {\em timeout} o memoria agotada, garantizando que los estadísticos describan el tiempo real necesario para demostrar optimalidad o hallar una solución heurística superior.

La robustez general de cada arquitectura se cuantificó mediante la razón incremental entre soluciones exitosas y el total de casos proyectados para cada entorno de simulación. La siguiente tabla sintetiza los resultados:

\blank[medium]

\bTABLE[split=yes, align=middle, frame=on, width=fit]
\setupTABLE[row][1][background=color, backgroundcolor=darkgray, style=bold]
\bTR \bTH Modelo \eTH \bTH Escenario \eTH \bTH Éxito (OK/Tot) \eTH \bTH Media (s) \eTH \bTH Desvío Est. (s) \eTH \bTH Máximo (s) \eTH \eTR
\bTR \bTD Exacto (OPP\_v3) \eTD \bTD Normal \eTD \bTD 48 / 48 \eTD \bTD 0.11 \eTD \bTD 0.14 \eTD \bTD 0.89 \eTD \eTR
\bTR \bTD Exacto (OPP\_v3) \eTD \bTD Stress \eTD \bTD 17 / 34 \eTD \bTD 164.45 \eTD \bTD 290.70 \eTD \bTD 912.54 \eTD \eTR
\bTR \bTD Hill Climbing \eTD \bTD Normal \eTD \bTD 46 / 46 \eTD \bTD 0.01 \eTD \bTD 0.00 \eTD \bTD 0.01 \eTD \eTR
\bTR \bTD Hill Climbing \eTD \bTD Stress \eTD \bTD 29 / 34 \eTD \bTD 10.41 \eTD \bTD 33.29 \eTD \bTD 177.48 \eTD \eTR
\bTR \bTD ACO \eTD \bTD Normal \eTD \bTD 48 / 48 \eTD \bTD 0.15 \eTD \bTD 0.21 \eTD \bTD 0.81 \eTD \eTR
\bTR \bTD ACO \eTD \bTD Stress \eTD \bTD 14 / 34 \eTD \bTD 34.75 \eTD \bTD 118.53 \eTD \bTD 596.39 \eTD \eTR
\eTABLE

\blank[medium]

A partir de la distribución estadística, emergen observaciones críticas para el diseño de soluciones de {\em routing}. Al incorporar el consolidado completo de estrés para el modelo OPP\_v3, se constata que logró resolver el 50\% de la muestra masiva (17 de 34 instancias) de manera exitosa dentro del límite de tiempo. No obstante, en estas resoluciones óptimas, la media temporal escala drásticamente a 164.45 segundos, con una extrema dispersión ($\sigma \approx 290.7$ s) y múltiples escenarios que consumieron cerca de 912 segundos. Esto ratifica que la optimización exacta no provee determinismo operativo estable en layouts estresados.

En contraste, Hill Climbing ratifica su posición dominante en ambientes de escala industrial. Mantiene una tasa de convergencia del 85\% en pruebas de estrés y ancla su tiempo de cómputo en un sobresaliente promedio de $\approx 10.4$ segundos, reduciendo la penalidad temporal por un factor de 15x respecto al solver lineal.

\stopsection

\startsection[title={2. Análisis de Correlación Paramétrica y Sensibilidad Algorítmica}]

Para aislar matemáticamente las causas subyacentes de las penalidades temporales, se computó el Coeficiente de Correlación de Pearson ($r$) evaluando la dependencia lineal entre el tiempo de convergencia exitosa y cada variable estructural de las instancias del LPRP.

\placeformula[-]
\startformula
r = \frac{\sum_{i=1}^{k} (x_i - \bar{x})(y_i - \bar{y})}{\sqrt{\sum_{i=1}^{k} (x_i - \bar{x})^2 \sum_{i=1}^{k} (y_i - \bar{y})^2}}
\stopformula

\blank[medium]

\bTABLE[split=yes, align=middle, frame=on, width=broad]
\setupTABLE[row][1][background=color, backgroundcolor=darkgray, style=bold]
\bTR \bTH Atributo del Problema ($x$) \eTH \bTH Exacto MILP ($r$) \eTH \bTH Hill Climbing ($r$) \eTH \bTH Metaheurística ACO ($r$) \eTH \eTR
\bTR \bTD Cantidad de Pedidos ($N$) \eTD \bTD 0.6597 \eTD \bTD 0.5347 \eTD \bTD 0.1802 \eTD \eTR
\bTR \bTD Vértices del Almacén ($V$) \eTD \bTD 0.2029 \eTD \bTD 0.1192 \eTD \bTD 0.9463 \eTD \eTR
\bTR \bTD Cantidad de Batches ($B$) \eTD \bTD 0.7085 \eTD \bTD 0.4963 \eTD \bTD 0.1327 \eTD \eTR
\bTR \bTD Capacidad del Vehículo ($C_v$) \eTD \bTD 0.4321 \eTD \bTD 0.3643 \eTD \bTD 0.0384 \eTD \eTR
\bTR \bTD Peso de Utilización ($W_{full}$) \eTD \bTD 0.3649 \eTD \bTD 0.2561 \eTD \bTD 0.3003 \eTD \eTR
\bTR \bTD Pen. Ventanas ($W_{ea\_ta}$) \eTD \bTD -0.2545 \eTD \bTD -0.1791 \eTD \bTD -0.1833 \eTD \eTR
\eTABLE

\blank[medium]

\startsubsection[title={2.1. Explosión Combinatoria del Batching frente al Volumen de Órdenes ($N$)}]

Al analizar la corrida completa, el modelo lineal OPP\_v3 manifiesta una sólida correlación ($r \approx 0.66$) respecto a los pedidos ($N$) y una correlación preponderante respecto a la cantidad de lotes o agrupaciones requeridas ($B \rightarrow r \approx 0.71$). La formulación del {\em order batching} introduce hiperdensidad en las matrices de asignación binaria ($z_{o,b}$). SCIP requiere expandir un árbol de ramificación y acotación exponencial que castiga el desempeño del solver, empujando la resolución de lotes de $>300$ pedidos hacia el límite del fracaso computacional.

Hill Climbing amortigua este castigo ($r \approx 0.53$) limitando la evaluación a vecindades con operadores polinomiales. ACO resulta prácticamente indiferente al empaquetamiento directo ($r \approx 0.18$) gracias a su construcción secuencial probabilística de soluciones.

\placefigure[here][fig:n_vs_time]{Curvas de Crecimiento Temporal vs Cantidad de Pedidos ($N$)}
{\startMPcode
  numeric u, v;
  u := 10cm / 500;
  v := 5cm / 1000;

  drawarrow (0,0) -- (550*u,0) withpen pencircle scaled 1pt;
  drawarrow (0,0) -- (0,1050*v) withpen pencircle scaled 1pt;
  
  label.bot("Cantidad de Pedidos (N)", (250*u, -40*v));
  label.lft("Tiempo Total (s)", (-15*u, 500*v));

  for i=0 step 200 until 1000:
    draw (-2*u, i*v) -- (2*u, i*v);
    label.lft(decimal i, (-5*u, i*v));
  endfor;

  for i=0 step 100 until 500:
    draw (i*u, -2*v) -- (i*u, 2*v);
    label.bot(decimal i, (i*u, -10*v));
  endfor;

  % Curva Exacto (Roja) ajustada a complete_output_v3_stress
  draw (12*u, 3.8*v) {dir 12} .. (100*u, 99.4*v) .. (320*u, 912.5*v) 
       withcolor \MPcolor{redline} withpen pencircle scaled 1.5pt;
  label.top("OPP Exacto (SCIP)", (200*u, 650*v)) withcolor \MPcolor{redline};

  % Curva Hill Climbing (Azul)
  draw (12*u, 0.01*v) -- (100*u, 23.3*v) -- (300*u, 80*v) -- (500*u, 177.4*v) 
       withcolor \MPcolor{blueline} withpen pencircle scaled 1.5pt;
  label.bot("Hill Climbing", (450*u, 230*v)) withcolor \MPcolor{blueline};
\stopMPcode}

\stopsubsection

\startsubsection[title={2.2. Colapso del Espacio de Estados en Metaheurísticas por Escala Topológica ($V$)}]

ACO sufre un colapso estructural al crecer la cantidad de vértices físicos en el almacén ($V$), registrando la máxima dependencia detectada en la tesis: $r = 0.9463$. Para escalar el layout a entornos masivos (625 nodos), la matriz probabilística de feromonas obliga a recalcular adyacencias iterativamente. Este impacto geométrico estanca asintóticamente al algoritmo frente al solver lineal, que ignora fuertemente la dimensión espacial ($r = 0.2029$ en la muestra global) favoreciendo las restricciones lógicas.

Hill Climbing brilla en esta faceta ($r = 0.1192$) al aislar computacionalmente el enrutamiento mediante grafos subyacentes precalculados, demostrando ser agnóstico al plano topológico.

\placefigure[here][fig:v_vs_time]{Curvas de Crecimiento Temporal vs Dimensión de Layout ($V$)}
{\startMPcode
  numeric u, v;
  u := 10cm / 600;
  v := 5cm / 600;

  drawarrow (0,0) -- (650*u,0) withpen pencircle scaled 1pt;
  drawarrow (0,0) -- (0,650*v) withpen pencircle scaled 1pt;
  
  label.bot("Vertices / Escala Topologica (V)", (300*u, -30*v));
  label.lft("Tiempo Total (s)", (-10*u, 300*v));

  for i=0 step 100 until 600:
    draw (-2*u, i*v) -- (2*u, i*v);
    label.lft(decimal i, (-3*u, i*v));
  endfor;

  for i=0 step 100 until 600:
    draw (i*u, -2*v) -- (i*u, 2*v);
    label.bot(decimal i, (i*u, -10*v));
  endfor;

  % Curva ACO (Verde)
  draw (16*u, 0.05*v) {dir 15} .. (100*u, 60.3*v) .. (400*u, 596.3*v) 
       withcolor \MPcolor{greenline} withpen pencircle scaled 1.5pt;
  label.lft("Metaheuristica ACO", (380*u, 540*v)) withcolor \MPcolor{greenline};

  % Curva Exacto en V (Roja) - Para la corrida completa tiene bajo impacto frente a N, pero es r=0.20
  draw (16*u, 0.07*v) -- (100*u, 36.1*v) -- (400*u, 92.8*v) -- (600*u, 243.1*v) 
       withcolor \MPcolor{redline} withpen pencircle scaled 1.5pt;
  label.ulft("OPP Exacto", (600*u, 260*v)) withcolor \MPcolor{redline};

  % Curva Hill Climbing (Azul)
  draw (16*u, 0.01*v) -- (100*u, 3*v) -- (400*u, 6*v) -- (600*u, 15*v) 
       withcolor \MPcolor{blueline} withpen pencircle scaled 1.5pt;
  label.bot("Hill Climbing", (500*u, 45*v)) withcolor \MPcolor{blueline};
\stopMPcode}

\stopsubsection

\startsubsection[title={2.3. Dinámica Restrictiva de la Capacidad Vehicular ($C_v$) y Densidad de Poda}]

La capacidad del vehículo evidencia que el empaque holgado corrompe al modelo exacto ($r = 0.4321$). Camiones sumamente restrictivos facilitan al solver SCIP aplicar cortes rápidos en ramas infactibles ({\em Branch and Cut}). Cuando se relaja $C_v$, SCIP debe validar un volumen abismal de soluciones hiperdensas. Las heurísticas HC y ACO neutralizan este efecto matemático mediante iteraciones secuenciales y asignaciones {\em greedy}, diluyendo su peso.

\stopsubsection

\startsubsection[title={2.4. Gradientes Lógicos y Penalizaciones Temporales ($W_{ea\_ta}$)}]

Resulta destacable que la variable de penalidad por ventanas de tiempo de entrega se comporte en régimen negativo para todas las plataformas (hasta $r = -0.2545$ en SCIP). Explicado algebraicamente, el endurecimiento de la función de costo objetivo castiga severamente el valor residual de rutas de ruteo ineficientes, expandiendo drásticamente la brecha con la Cota Dual y permitiendo descartar escenarios tempranamente. Castigar retrasos optimiza la eficiencia del sistema.

\stopsubsection

\stopsection

\startsection[title={3. Conclusiones Fundacionales y Directrices de Diseño Logístico}]

\startitemize[n]
\item {\bf Escala Límite del Modelo MILP:} La consolidación de las bitácoras completas ratifica la inhabilidad de OPP\_v3 para brindar un desempeño garantizado en tiempos industriales bajo grandes volúmenes de pedidos o hiper-lotes. Sus tasas de éxito fluctúan violentamente al incorporar ruido en la variable ($N$), tornándolo inseguro para sistemas dinámicos (WMS).
\item {\bf Sensibilidad Geográfica de ACO:} La exploración algorítmica expone que ACO no sufre por qué empaquetar, sino por {\em dónde}. Al escalar la red topológica, el estancamiento computacional provocado por el denso recálculo matricial de grafos desaconseja su uso sin aplicar heurísticas severas de segmentación de zonas de picking.
\item {\bf Dominancia Heurística Consolidada:} Hill Climbing valida su rol de arquitectura cúspide frente al conjunto completo de pruebas. Inmune a la dimensión del plano geográfico gracias al aislamiento en su matriz de distancias y resistente al incremento de la densidad de órdenes, asegura una convergencia excepcionalmente ágil ($\mu = 10.41$ s en estrés) y establece la mejor relación de compromiso factibilidad/tiempo.
\stopitemize

\stopsection

\stoptext
